% Единственный источник содержания темы. UTF-8.
% Метки label и ключи в [...] сохраняют опубликованные ссылки.

\section{Теория}

\subsection{Определение}

\begin{task}[conditional-probability/theory/1]\label{task-theory-1}
\textbf{Задача 1.} а) В семье имеется два ребёнка (не близнецы и не двойняшки). Какова вероятность, что оба ребёнка  —  девочки?\\
 б) В семье имеется два ребёнка (не близнецы и не двойняшки), один из которых девочка. Какова вероятность, что оба ребёнка  —  девочки?
\end{task}

\begin{solution}[conditional-probability/theory/1]\label{solution-conditional-probability-theory-1-method-1}
\textbf{Решение.} а) Пространство элементарных исходов $\Omega = \{\text{ММ, ДД, МД, ДМ}\}$ (так как есть старший и младший, мы отличаем МД и ДМ). Тогда у нас всего одно благоприятсвующее элементарное событие и вероятность равна $\frac{1}{4}$.\\
 б) Мы уже знаем, что один из детей  —  девочка, значит исхода ММ у нас нет (или же можно считать что он есть, но невозможен), и $\Omega = \{ \text{МД, ДМ, ДД}\}$. Вероятность соответственно равна $\frac{1}{3}$.\\
\end{solution}

Чтобы решить пункт б) нам пришлось изменять пространство элементарных исходов, что часто бывает неудобно. Например, чтобы вычислить вероятность того, что второй ребёнок мальчик в семье с двумя детьми, один из которых девочка, вновь придётся менять вероятностное пространство. Чтобы избежать этого, мы будем использовать \emph{условную вероятность}.

\begin{designation}
Вероятность наступления события $A$ при условии того, что наступило событие $H$ обозначаем $\mathbb{P}(A\vert H)$
\end{designation}

\begin{definition}
\textbf{Условная вероятность}. Зафиксируем некоторое событие $H$, такое, что $\mathbb{P}(H) > 0$. Тогда вероятность наступления события $A$ при условии, что событие $H$ наступило равна

\[\mathbb{P}(A\vert H) = \frac{\mathbb{P}(A\cap H)}{\mathbb{P}(H)}\]
\end{definition}

\begin{solution}[conditional-probability/theory/1]\label{solution-conditional-probability-theory-1-conditional-method}
\textbf{Решение задачи 1. \emph{Условная вероятность}} Пусть $A = \{$ В семье обе девочки $\} = \{"\text{ДД}"\}$ и $H = \{\text{в семье есть одна девочка}\} = \{\text{"МД", "ДМ", "ДД"}\}$. Тогда $A\cap H = \{\text{ДД}\}$ и $\mathbb{P}(A\vert H) = \frac{\mathbb{P}(A\cap H)}{\mathbb{P}(H)} = \frac{1/4}{3/4} = \frac{1}{3}$.
\end{solution}

Иногда наоборот мы можем легко вычислить условную вероятность и нам требуется вычислить вероятность наступления обоих событий, то есть вероятность пересечения. Тогда удобно переписать последнюю формулу в следующем виде (работает только если каждое из событий $A, B$ не с нулевой вероятностью).

\begin{theorem}[conditional-probability/theorem/1]\label{theorem-1}
\textbf{Формула умножения вероятностей.} $\mathbb{P}(A\cap B) = \mathbb{P}(B)\cdot \mathbb{P}(A\vert B) = \mathbb{P}(A)\cdot \mathbb{P}(B\vert A)$
\end{theorem}

\begin{proof}[conditional-probability/theorem/1]\label{solution-conditional-probability-theorem-1}
При $\mathbb P(B)>0$ определение условной вероятности даёт

\[\mathbb P(A\mid B)=\frac{\mathbb P(A\cap B)}{\mathbb P(B)}.\]

Умножив на $\mathbb P(B)$, получаем $\mathbb P(A\cap B)=\mathbb P(B)\mathbb P(A\mid B)$.

Если также $\mathbb P(A)>0$, можно поменять события местами и получить $\mathbb P(A\cap B)=\mathbb P(A)\mathbb P(B\mid A)$. При положительных вероятностях обоих событий обе записи определены и равны между собой.
\end{proof}

\begin{task}[conditional-probability/theory/2]\label{task-theory-2}
\textbf{Задача 2.} В игре под одноимённым названием из фильма "Жмурки" два игрока используют шестизарядные револьвер с одним патроном. Первый раскручивает барабан револьвера и нажимает на спусковой крючок, затем аналогичное действие выполняет второй, будем считать что дальше игра заканчивается. Найдите вероятность того, что а) оба нажатия не произведут выстрел, б) выстрел произойдёт при втором нажатии.
\end{task}

\begin{solution}[conditional-probability/theory/2]\label{solution-conditional-probability-theory-2-method-1}
\textbf{Решение.} а) Пусть $A = \{\text{Первый не выстрелил}\}$ и $B =\{\text{Второй не выстрелил}\}$. Тогда $\mathbb{P}(B\vert A) = 5/6$ так как второй раскручивает барабан (тем самым "сбрасывая вероятность") в котором всё ещё один патрон. И по формуле умножения вероятностей $\mathbb{P}(A\cap B)= \mathbb{P}(A)\cdot \mathbb{P}(B|A) = 5/6 \cdot 5/6 = \frac{25}{36}$. Если бы второй не раскручивал барабан $\mathbb{P}(B\vert A) = 4/5$ так как при его ходе осталось 5 мест в барабане одно из которых всё ещё с патроном. И соответственно $\mathbb{P}(A\cap B) = 4/6 = 2/3$, то есть вероятность в таком случае меньше, и вероятность, что кто то выстрелит соответственно больше.\\
 б) Теперь пусть $C = \{\text{Второй выстрелил}\}$ тогда требуется найти $\mathbb{P}(C\cap A) = \mathbb{P}(C\vert A)\cdot \mathbb{P}(A) = 5/6 \cdot 1/6 = 5/36$ и $5/6\cdot 1/5 = 1/6$, если второй не раскручивает барабан.
\end{solution}

\section{Задачи}

\begin{task}[conditional-probability/problems/1]\label{task-problems-1}
\textbf{3.} В торговом центре установлены два автомата, продающие кофе. Вероятность того, что к концу дня кофе закончится в каждом отдельном автомате, равна 0,3. В обоих автоматах кофе заканчивается к вечеру с вероятностью 0,21. Вечером пришёл мастер, чтобы обслужить автоматы, и обнаружил, что во втором автомате кофе закончился. Какова теперь вероятность, что и в первом автомате уже нет кофе?
\end{task}

\begin{solution}[conditional-probability/problems/1]\label{solution-conditional-probability-problems-1-method-1}
\textbf{Решение.} Пусть $A=$ \{"в первом автомате кончился кофе "\} и $B$ = \{"во втором закончился кофе\}. Тогда $\mathbb{P}(A\vert B) = \frac{\mathbb{P}(A\cap B)}{\mathbb{P}(B)} = \frac{0,21}{0,3} = 0,7$.
\end{solution}

\begin{task}[conditional-probability/problems/2]\label{task-problems-2}
\textbf{4.} В некотором городе 7\% населения — студенты. Из всех студентов 60\% учатся в университете. Найдите вероятность того, что случайно выбранный житель этого города является студентом университета.
\end{task}

\begin{solution}[conditional-probability/problems/2]\label{solution-conditional-probability-problems-2-method-1}
\textbf{Решение 1.} Пусть в городе $N$ жителей, тогда студентов университета $N\cdot \frac{7}{100}\cdot \frac{60}{100} = \frac{42}{1000}\cdot N$. Это количество студентов университета, то есть количество благоприятствующих событий. Вероятность получается делением этого выражения на $N$\\
\end{solution}

\begin{solution}[conditional-probability/problems/2]\label{solution-conditional-probability-problems-2-method-2}
\textbf{Решение 2.} Пусть $A$ = \{"Житель  —  студент"\}, $B=$ \{"Житель  —  учится в университете"\}. Тогда $\mathbb{P}(A\cap B) = \mathbb{P}(A)\cdot \mathbb{P}(B\vert A) = 0,07 \cdot 0,6=0,042$.
\end{solution}

\begin{task}[conditional-probability/problems/3]\label{task-problems-3}
\textbf{5.} Найдите вероятность получить $5$ разных результатов, при $5$ подбрасываниях игрального кубика.
\end{task}

\begin{solution}[conditional-probability/problems/3]\label{solution-conditional-probability-problems-3-method-1}
\textbf{Решение.} Пусть $A_k$ = \{"$k$-ый бросок дал результат отличающийся от каждого из предыдущих результатов"\}, где $k \in \{1, 2,3,4,5\}$ (считаем, что результат первого броска отличается от результата нуля бросков, так как первый бросок уникальный). Тогда $\mathbb{P}(A_2|A_1)$  —  вероятность того что первые два броска различны. $\mathbb{P}(A_3\vert A_2\cap A_1)$  —  вероятность того что третий бросок отличается от первых двух, при условии что второй бросок отличается от первого, то есть что все три результата различны. Аналогично продолжаем и получаем что нам требуется вычислить $\mathbb{P}(A_5\cap A_4\cap A_3\cap A_2\cap A_1)$. Теперь будем последовательно считать вероятности.\\
 1)$\mathbb{P}(A_2| A_1) = 5/6$, так как первый бросок всегда уникальный, а второй будет отличаться от него, если и только если результат не попал в одно единственное число.\\
 2) $p_3=\mathbb{P}(A_3\cap A_2\cap A_1) = \mathbb{P}(A_3\cap (A_2\cap A_1)) = \mathbb{P}(A_3\vert A_2\cap A_1)\cdot \mathbb{P}(A_2\cap A_1) = 4/6\cdot 5/6 = \frac{4\cdot 5}{6^2}$.\\
 3) $p_4 = \mathbb{P}(A_4\cap A_3 \cap A_2\cap A_1) = \mathbb{P}(A_4\cap (A_3 \cap A_2\cap A_1)) = \mathbb{P}(A_4\vert A_3\cap A_2\cap A_1)\cdot p_3 = 3/6 \cdot p_3 = \frac{3\cdot 4\cdot 5}{6^3}$.\\
 4) $p_5 = \mathbb{P}(A_5\cap A_4\cap A_3\cap A_2\cap A_1) = \mathbb{P}(A_5\cap (A_4\cap A_3\cap A_2\cap A_1)) = \mathbb{P}(A_5\vert  A_4\cap A_3\cap A_2\cap A_1)\cdot p_4 = p_4\cdot 2/6 = \frac{2\cdot 3\cdot 4\cdot 5}{6^4} = \frac{4\cdot 5}{6^3} = \frac{5}{3^3\cdot 2}=\frac{5}{54}$.
\end{solution}

\section{Домашнее задание}

Вторая часть, стр. 36--39

\begin{task}[conditional-probability/homework/1]\label{task-homework-1}
\textbf{1. №101} При двукратном бросании монеты в первый раз выпала решка. Найдите условную вероятность события:\\
 а) «оба раза выпадет решка»;\\
 б) «выпадет хотя бы один орёл»;\\
 в) «выпадут два орла».
\end{task}

\begin{solution}[conditional-probability/homework/1]\label{solution-conditional-probability-homework-1}
Из четырёх равновозможных исходов двух бросков ОО, ОР, РО, РР условию «первая — решка» соответствуют только РО и РР. После этого условия они остаются равновозможными.

а) Две решки — исход РР: вероятность $\frac12$.

б) Хотя бы один орёл — исход РО: вероятность $\frac12$.

в) Два орла противоречат известному результату первого броска: вероятность 0.
\end{solution}

\begin{task}[conditional-probability/homework/2]\label{task-homework-2}
\textbf{2. №103} При двукратном бросании игральной кости сумма выпавших очков равна 9. Найдите условную вероятность события:\\
 а) «в первый раз выпадет 5 очков»;\\
 б) «при одном из бросков выпадет 4 очка»;\\
 в) «в первый раз выпадет меньше очков, чем во второй»;\\
 г) «во второй раз выпадет меньше чем 3 очка».
\end{task}

\begin{solution}[conditional-probability/homework/2]\label{solution-conditional-probability-homework-2}
При сумме 9 остаются четыре равновозможных упорядоченных исхода:

\[H=\{(3,6),(4,5),(5,4),(6,3)\}.\]

а) Первые 5 очков — только исход $(5,4)$: вероятность $\frac14$.

б) Четыре очка при одном из бросков — исходы $(4,5),(5,4)$: вероятность $\frac24=\frac12$.

в) В первый раз меньше, чем во второй, — исходы $(3,6),(4,5)$: вероятность $\frac12$.

г) Во второй раз меньше 3 очков — подходящих исходов нет: вероятность 0.
\end{solution}

\begin{task}[conditional-probability/homework/3]\label{task-homework-3}
\textbf{3.} №107. В посёлке 40\% взрослого населения занято в сельском хозяйстве, причём 5\% взрослого населения посёлка работают в агропромышленном холдинге «Нива». Для опроса случайно выбран житель этого посёлка, и оказалось, что он занят в сельском хозяйстве. При этом условии найдите условную вероятность того, что он работает в холдинге «Нива».\\
\end{task}

\begin{solution}[conditional-probability/homework/3]\label{solution-conditional-probability-homework-3-method-1}
\textbf{Решение.} Пусть $A =$ \{"человек занят в С/Х"\} и $B =$ \{"человек работает в холдинге "Нива"\}. Видимо авторы учебника подразумевают, что все работники холдинга занимаются сельским хозяйством, что вполне логично. Тогда $\mathbb{P}(B|A)=\frac{\mathbb{P}(A\cap B)}{\mathbb{P}(A)}$, где $\mathbb{P}(A) = 0,4$, а $\mathbb{P}(A\cap B) = \mathbb{P}(B) = 0,05$. То есть $\mathbb{P}(B|A) = \frac{1}{8}$. Важно, что раз любой рабочий холдинга занят сельским хозяйством, то событие $B$ лежит в событии $A$, а значит они пересекаются по событию $B$.
\end{solution}

\begin{task}[conditional-probability/homework/4]\label{task-homework-4}
\textbf{4. №109 (г, д)} Найдите вероятность получить $n$ разных результатов, бросив игральную кость $n$ раз, если:\\
 г) $n=6$;\\
 д) $n=7$.
\end{task}

\begin{solution}[conditional-probability/homework/4]\label{solution-conditional-probability-homework-4}
г) Чтобы шесть результатов оказались различными, каждый очередной бросок должен дать ещё не выпадавшее число. По формуле умножения вероятностей

\[\mathbb P=\frac66\cdot\frac56\cdot\frac46\cdot\frac36\cdot\frac26\cdot\frac16=\frac{720}{46656}=\frac5{324}.\]

д) У кубика только шесть возможных результатов. Получить семь разных результатов невозможно, поэтому вероятность равна 0.
\end{solution}

\begin{task}[conditional-probability/homework/5]\label{task-homework-5}
\textbf{5.} №110. В ящике 20 левых и 20 правых перчаток. Сколько нужно достать перчаток, не глядя в ящик, чтобы среди вынутых перчаток нашлась хотя бы одна левая и хотя бы одна правая перчатка:\\
 б) с вероятностью не меньше чем $0,95$.\\
\end{task}

\begin{solution}[conditional-probability/homework/5]\label{solution-conditional-probability-homework-5-method-1}
\textbf{Решение пункта б.} \\
 Так как мы хотим узнать количество перчаток, которые требуется достать, обозначим это количество за неизвестную $n$. Пусть мы достали $x$ правых перчаток и $y$  —  левых, то есть $x+y = n$. Пусть $A$  —  событие противоположное требуемому, то есть $A =$ \{"среди вытянутых перчаток 0 левых перчаток или 0 правых\}. Тогда пусть $L =$ \{"Среди вытянутых перчаток 0 левых"\} и $R =$ \{"Среди вытянутых перчаток 0 правых"\}. Отсюда $A= L\cup R$. Чтобы получить $\mathbb{P}(\overline{A}) \geq 0,95$, достаточно потребовать, чтобы $\mathbb{P}(A) = 1 - \mathbb{P}(\overline{A}) < 0,05$ или как мы написали $\mathbb{P}(L\cup R) < 0,05$. По формуле сложения вероятностей $\mathbb{P}(L\cup R) = \mathbb{P}(L) + \mathbb{P}(R) - \mathbb{P}(L\cap R)$. Ясно, что $n= 1$ нам не подходит, а значит при других $n>1$ нельзя одновременно не вытащить ни левую, ни правую, то есть $\mathbb{P}(L\cap R) = 0$. И мы хотим чтобы $\mathbb{P}(L) + \mathbb{P}(R) < 0,05$. Теперь посчитаем, $\mathbb{P}(R)$ как функцию от $n$. Мы хотим вытащить все $n$ левых перчаток из $40$, тогда по аналогии с задачей \textbf{5} урока, вероятность вытащить первой правую перчатку равна $\frac{20}{40}$, вероятность, что вторая будет тоже правая равна $\frac{19}{39}$ (осталось 19 правых и 39 всего), вероятность что третья тоже правая $\frac{18}{38}$ и так далее до $n$-ой перчатки с вероятностью $\frac{20-n +1}{40 - n +1}$. Итого $\mathbb{P}(R) = \frac{20}{40}\cdot \frac{19}{39}\cdot\frac{18}{38}\cdot \dots \frac{21 - n}{41-n}$. Так правые и левые перчатки симметричны $\mathbb{P}(L) = \mathbb{P}(R) = \frac{20}{40}\cdot \frac{19}{39}\cdot\frac{18}{38}\cdot \dots \frac{21 - n}{41-n}$ и $\mathbb{P}(A) = 2\cdot \frac{20\cdot 19\cdot\dots\cdot (21-n)}{40\cdot 39\cdot \dots\cdot (41 - n)}$. Остаётся решить неравенство $2\cdot \frac{20\cdot 19\cdot\dots\cdot (21-n)}{40\cdot 39\cdot \dots\cdot (41 - n)} < 0,05$. Проще всего сделать это перебором различных $n$ от $2$ по возрастанию и считать результаты на калькуляторе. Других способов без калькулятора я не знаю и не думаю, что они есть. При $n=4$ вероятность неудачи равна $\frac{51}{481}>0,05$, а при $n=5$ — $\frac{68}{1443}<0,05$. При меньшем числе перчаток вероятность неудачи больше. Следовательно, наименьшее подходящее количество перчаток — $n=5$.
\end{solution}

